% !TeX root = ../../Book1.tex
\section{一致可积性与积分收敛}
\label{b1:sec:3106}

上一节得到的极限仍留下两个问题：何时能让极限进入积分，
又何时能把每个有限时刻的鞅值解释为对最终值的条件平均？
这两件事由同一个条件控制。
\secref{b1:sec:0505}的一致可积性排除积分质量在小概率的大值事件上集中，
而鞅结构把 \(L^1\) 收敛进一步转成终值表示。
本节先证明这个联系，再用它识别随信息增加而形成的条件平均列的极限。

\subsection{一致可积鞅}
\label{b1:subsec:310601}

在概率空间 \((\Omega,\mathcal F,\mathbb P)\) 上，实值或复值可积随机变量族
\(\mathcal K\) 一致可积，指
\[
 \lim_{R\to\infty}\sup_{X\in\mathcal K}
       \mathbb E\bigl[|X|\cdot\mathbf1_{\{|X|>R\}}\bigr]=0.
\]
这就是\cref{b1:def:uniform-integrability}在总质量为一的空间上的定义。
称一列鞅一致可积，是要求它的全体有限时刻值构成这样一个族。
阈值 \(R\) 必须对所有时刻同时有效；对每个固定时刻分别取一个阈值，只是通常的可积性。

\ProseParagraphBreak

因为全空间概率有限，\cref{b1:prop:uniform-integrability-criterion}说明，
一致可积性等价于一阶矩一致有界，加上积分对小概率事件的一致绝对连续性。
具体说，后一条件要求：给定 \(\varepsilon>0\)，存在 \(\delta>0\)，
使 \(\mathbb P(A)<\delta\) 时，\(\sup_{X\in\mathcal K}\mathbb E[|X|\cdot\mathbf1_A]<\varepsilon\)。
因此，\eqref{b1:eq:dyadic-spike-martingale-31}的失败正好落在这个缺失的条件上。

\begin{lemma}[条件平均的一致可积性]
\label{b1:lem:conditional-expectations-ui-31}
设 \((\Omega,\mathcal F,\mathbb P)\) 为概率空间，\(Y\in L^1(\mathbb P)\)
为实值或复值随机变量。让 \(\mathcal G\) 遍历 \(\mathcal F\) 的任意一族子
\(\sigma\)-代数，则所有 \(\mathbb E[Y\mid\mathcal G]\) 构成一致可积族。
更精确地，对每个 \(K>0\)，令 \(Y_{\mathcal G}=\mathbb E[Y\mid\mathcal G]\)，则
\begin{equation}
 \mathbb E\bigl[|Y_{\mathcal G}|\cdot
                 \mathbf1_{\{|Y_{\mathcal G}|>K\}}\bigr]
 \leq2\mathbb E\bigl[|Y|\cdot\mathbf1_{\{|Y|>K/2\}}\bigr].
 \label{b1:eq:conditional-ui-tail-31}
\end{equation}
\end{lemma}
\begin{proof}
由条件期望的绝对值估计和正性，
\[
 |Y_{\mathcal G}|
 \leq\mathbb E[|Y|\mid\mathcal G]
 \leq K/2+\mathbb E\bigl[|Y|\cdot\mathbf1_{\{|Y|>K/2\}}\mid\mathcal G\bigr].
\]
集合 \(A=\{|Y_{\mathcal G}|>K\}\) 属于 \(\mathcal G\)。在这个集合上，
\(K/2<|Y_{\mathcal G}|/2\)，故
\[
 |Y_{\mathcal G}|\cdot\mathbf1_A
 \leq2\mathbf1_A\cdot
 \mathbb E\bigl[|Y|\cdot\mathbf1_{\{|Y|>K/2\}}\mid\mathcal G\bigr].
\]
对 \(A\) 使用条件期望的积分恒等式，再把非负被积函数的积分扩大到全空间，
就得到\eqref{b1:eq:conditional-ui-tail-31}。
右侧与 \(\mathcal G\) 无关，并由 \(Y\in L^1\) 在 \(K\to\infty\) 时趋零。
\end{proof}

这个结论不要求子 \(\sigma\)-代数之间相互包含。
对固定的一个可积变量取任何信息下的平均，都不会产生无法统一控制的积分尖峰。
注意这里固定的是 \(Y\)；若把 \(Y\) 也任意改变，只有 \(L^1\) 有界性就不足以得到结论。

\begin{proposition}[较高阶矩给出一致可积性]
\label{b1:prop:lp-ui-31}
设 \((\Omega,\mathcal F,\mathbb P)\) 为概率空间，\(\mathcal K\) 为其中实值或复值随机变量族。
若某个 \(1<p<\infty\) 满足 \(\sup_{X\in\mathcal K}\mathbb E|X|^p=C<\infty\)，
则 \(\mathcal K\subset L^1\) 且一致可积，并有
\[
 \sup_{X\in\mathcal K}\mathbb E\bigl[|X|\cdot\mathbf1_{\{|X|>R\}}\bigr]
 \leq C R^{1-p},\quad R>0.
\]
若全族的 \(L^\infty\) 范数有共同有限上界，也一致可积。
\end{proposition}
\begin{proof}
由概率空间上的 Hölder 不等式，\(\|X\|_1\leq\|X\|_p\)。
在 \(\{|X|>R\}\) 上，\(|X|\leq R^{1-p}|X|^p\)，积分即得估计，
且 \(1-p<0\)。本质有界的情形在 \(R\) 超过共同上界时，所有尾部积分均为零。
\end{proof}

\subsection{闭鞅与终值表示}
\label{b1:subsec:310602}

有限时间区间上的鞅总满足 \(M_n=\mathbb E[M_N\mid\mathcal F_n]\)，\(n\leq N\)。
在无限时间区间上，是否仍能找到一个可积的共同终值？
\eqref{b1:eq:dyadic-spike-martingale-31}已经说明答案并非总是肯定的。
下面给出准确判据。

\begin{definition}[闭鞅]
\label{b1:def:closed-martingale-31}
设 \((\Omega,\mathcal F,\mathbb P)\) 为概率空间，\((\mathcal F_n)_{n\geq0}\)
为递增子 \(\sigma\)-代数列，\((M_n)\) 为适应于该滤过的实鞅或复鞅。
若存在 \(Y\in L^1(\Omega,\mathcal F,\mathbb P)\)，使每个 \(n\geq0\) 都满足
\[
 M_n=\mathbb E[Y\mid\mathcal F_n]\quad\text{几乎必然},
\]
就称 \((M_n)\) 是\term[biyang]{闭鞅}[closed martingale]，或称它由 \(Y\) 闭合。
\end{definition}

闭合变量 \(Y\) 不要求已经对 \(\mathcal F_\infty=\sigma(\bigcup_n\mathcal F_n)\) 可测。
例如滤过始终只有平凡信息，则 \(M_n=\mathbb E Y\) 对所有 \(n\) 成立，
而 \(Y\) 可以是任何具有该均值的可积变量。因此共同终值表示中的 \(Y\) 未必唯一。
真正由这一滤过看见的终值将在下一小节识别。

\begin{theorem}[一致可积鞅与终值表示]
\label{b1:thm:closed-martingale-31}
设 \((\Omega,\mathcal F,\mathbb P)\) 为概率空间，\((\mathcal F_n)_{n\geq0}\)
为递增子 \(\sigma\)-代数列，\((M_n)\) 为适应于该滤过的实鞅或复鞅。
下列三项等价：
\begin{equivalences}
 \item\label{b1:item:closed-ui-31}\(\{M_n:n\geq0\}\) 一致可积。
 \item\label{b1:item:closed-l1-31}存在 \(M_\infty\in L^1(\mathbb P)\)，使 \(M_n\to M_\infty\) 于 \(L^1\)。
 \item\label{b1:item:closed-representation-31}\((M_n)\) 是闭鞅。
\end{equivalences}
这些条件成立时，\(M_n\) 还几乎必然收敛于 \(M_\infty\)，
可取 \(M_\infty\) 对 \(\mathcal F_\infty=\sigma(\bigcup_n\mathcal F_n)\) 可测，且
\begin{equation}
 M_n=\mathbb E[M_\infty\mid\mathcal F_n],\quad n\geq0.
 \label{b1:eq:martingale-terminal-representation-31}
\end{equation}
特别地，\(\mathbb E M_\infty=\mathbb E M_0\)。
\end{theorem}
\begin{proof}
先设\itemref{b1:item:closed-ui-31}成立。概率空间上的一致可积性给出
\(\sup_n\mathbb E|M_n|<\infty\)，由\cref{b1:cor:martingale-as-convergence-31}，
存在可积的几乎必然极限 \(M_\infty\)。按有限极限存在的集合取代表，
可以使它 \(\mathcal F_\infty\)-可测。
有限测度空间上几乎处处收敛推出依测度收敛，因而\cref{b1:thm:vitali-convergence}
适用：全列一致可积，且依概率趋于 \(M_\infty\)，所以 \(\|M_n-M_\infty\|_1\to0\)。
这证明\itemref{b1:item:closed-l1-31}。

现设\itemref{b1:item:closed-l1-31}成立。固定 \(n\) 及 \(A\in\mathcal F_n\)。
对每个 \(m\geq n\)，鞅的多步条件平均关系给出
\[
 \mathbb E[M_m\cdot\mathbf1_A]=\mathbb E[M_n\cdot\mathbf1_A].
\]
左侧与 \(\mathbb E[M_\infty\cdot\mathbf1_A]\) 的差的绝对值不超过
\(\|M_m-M_\infty\|_1\)，因此可以令 \(m\to\infty\)。
又因 \(M_n\) 是 \(\mathcal F_n\)-可测的可积变量，条件期望的定义和唯一性便给出
\eqref{b1:eq:martingale-terminal-representation-31}。于是鞅由 \(M_\infty\) 闭合，得到\itemref{b1:item:closed-representation-31}。

最后，若\itemref{b1:item:closed-representation-31}成立，写
\(M_n=\mathbb E[Y\mid\mathcal F_n]\)，其中 \(Y\in L^1\)。
\cref{b1:lem:conditional-expectations-ui-31}直接给出\itemref{b1:item:closed-ui-31}。
于是也得到前面证明的几乎必然极限；它与原有 \(L^1\) 极限相等几乎必然，
因为 \(L^1\) 收敛推出依概率收敛，而依概率极限唯一。
取 \(A=\Omega\)，或直接在 \(L^1\) 极限中取期望，即得期望保持。
\end{proof}

Fremlin 的《Measure Theory》\cite[275H]{Fremlin2016Measure}
还把终值的绝对值积分与有限时刻积分的极限相等列作等价条件。
本书采用 \(L^1\) 收敛表述，因为它直接给出上面逐事件传递积分恒等式所需的误差估计。
% 实读：chap27 摘本物理 pp.52--53，275H陈述及证明。

\begin{corollary}[一致高阶矩下的积分收敛]
\label{b1:cor:lp-martingale-convergence-31}
设 \((M_n)_{n\geq0}\) 为概率空间 \((\Omega,\mathcal F,\mathbb P)\) 上
适应于递增滤过 \((\mathcal F_n)\) 的实鞅或复鞅。
若 \(1<p<\infty\) 且 \(\sup_n\|M_n\|_p<\infty\)，则有
\(M_\infty\in L^p\)，使 \(M_n\to M_\infty\) 几乎必然且于 \(L^p\)。
若 \(\sup_n\|M_n\|_\infty<\infty\)，则上述结论对每个有限 \(p\geq1\) 成立，
但不保证 \(L^\infty\) 收敛。
\end{corollary}
\begin{proof}
由\cref{b1:prop:lp-ui-31}及上一条定理，先得到几乎必然和 \(L^1\) 收敛。
Fatou 引理给出 \(M_\infty\in L^p\)。为了提升到 \(L^p\)，
还需要对 \(|M_n|^p\) 的积分控制，不能仅用其积分有界代替一致可积性。
对非负下鞅 \(|M_n|\) 使用\cref{b1:thm:doob-strong-31}，得
\[
 \mathbb E\max_{0\leq k\leq N}|M_k|^p
 \leq\left(\frac p{p-1}\right)^p\sup_n\mathbb E|M_n|^p.
\]
由单调收敛，\(M^*=\sup_{n\geq0}|M_n|\) 属于 \(L^p\)。
几乎处处又有 \(|M_\infty|\leq M^*\)，所以
\[
 |M_n-M_\infty|^p\leq 2^p(M^*)^p.
\]
右侧为与 \(n\) 无关的可积控制，控制收敛定理给出 \(L^p\) 收敛。
本质一致有界时，先取全部时刻例外零集的可数并；
在其补集上全列及极限都被同一个常数控制。
再对任意有限 \(p\) 应用控制收敛即可。
\end{proof}

\subsection{条件平均列的收敛}
\label{b1:subsec:310603}

设对同一个可积变量 \(Y\) 的信息逐步增加。
塔性质保证 \(\mathbb E[Y\mid\mathcal F_n]\) 构成鞅，上一条定理保证它有积分极限。
还须证明这个极限究竟是什么。
所有有限时刻能看见的事件的并一般不是 \(\sigma\)-代数；
因此从每个 \(\mathcal F_n\) 上的积分恒等式走到生成的 \(\mathcal F_\infty\)，需要一次明确的测度论延伸。

\begin{theorem}[Lévy 向上收敛]
\label{b1:thm:levy-upward-31}
设 \((\Omega,\mathcal F,\mathbb P)\) 为概率空间，\((\mathcal F_n)_{n\geq0}\)
为递增子 \(\sigma\)-代数列，并令
\(\mathcal F_\infty=\sigma(\bigcup_{n\geq0}\mathcal F_n)\)。
对每个实值或复值 \(Y\in L^1(\mathbb P)\)，都有
\[
 \mathbb E[Y\mid\mathcal F_n]
 \longrightarrow\mathbb E[Y\mid\mathcal F_\infty]
 \quad\text{几乎必然且于 }L^1.
\]
若 \(Y\) 有 \(\mathcal F_\infty\)-可测代表，极限就是 \(Y\)。
\end{theorem}
\begin{proof}
令 \(M_n=\mathbb E[Y\mid\mathcal F_n]\)。塔性质给出鞅关系，
而\cref{b1:lem:conditional-expectations-ui-31}给出一致可积性。
由\cref{b1:thm:closed-martingale-31}，存在
\(\mathcal F_\infty\)-可测的 \(M_\infty\in L^1\)，且 \(M_n\to M_\infty\) 几乎必然及于 \(L^1\)。
对 \(A\in\bigcup_n\mathcal F_n\)，可以选 \(j\) 使 \(A\in\mathcal F_j\)。
于是对所有 \(n\geq j\)，\(\mathbb E[M_n\cdot\mathbf1_A]=\mathbb E[Y\cdot\mathbf1_A]\)，
令 \(n\to\infty\) 得到相同的 \(M_\infty\) 积分恒等式。

记 \(\mathcal A_0=\bigcup_n\mathcal F_n\)。由于子代数递增，
任意有限多个 \(\mathcal A_0\) 中的集合同时属于某个 \(\mathcal F_j\)，
所以 \(\mathcal A_0\) 含 \(\Omega\) 且对有限交封闭。
再令
\[
 \mathcal L=\{\,A\in\mathcal F_\infty:
       \mathbb E[M_\infty\cdot\mathbf1_A]=\mathbb E[Y\cdot\mathbf1_A]\,\}.
\]
全空间属于 \(\mathcal L\)；从全空间积分减去 \(A\) 上的积分，说明它对取补封闭。
若 \(A_k\in\mathcal L\) 两两不交，则对前 \(m\) 个集合相加后令 \(m\to\infty\)。
被积函数受 \(|M_\infty|+|Y|\) 控制，故控制收敛保证并集也属于 \(\mathcal L\)。
因此 \(\mathcal L\) 是包含 \(\mathcal A_0\) 的 Dynkin \(\lambda\)-系统。
\cref{b1:thm:pi-lambda}给出 \(\mathcal L=\mathcal F_\infty\)。

此时 \(M_\infty\) 对 \(\mathcal F_\infty\) 可测、可积，
并在该子代数的每个事件上保持 \(Y\) 的积分，故由条件期望的唯一性，
\(M_\infty=\mathbb E[Y\mid\mathcal F_\infty]\)。
若 \(Y\) 有此子代数的可测代表，条件期望就是 \(Y\) 本身。
\end{proof}

这一\term[levy-xiangshang-shoulian-dingli]{Lévy 向上收敛定理}[Lévy's upward convergence theorem]
把“信息最终达到什么程度”精确表示为生成的 \(\sigma\)-代数。
这里的“向上”指信息递增，不指条件平均的数值单调。
对应的极限识别可对照 Fremlin 的《Measure Theory》\cite[275I]{Fremlin2016Measure}。
% 实读：chap27 摘本物理 p.53 的275I；不把后继275K反向鞅纳入本节已证范围。

\ProseParagraphBreak

再看区间上的二进信息。仍令 \(\mathcal D_n\) 由长度 \(2^{-n}\) 的半开区间
\(I_{n,j}\subset[0,1)\) 生成。对任意 Borel 可测的 \(f\in L^1([0,1))\)，令
\[
 A_nf(x)=2^n\int_{I_{n,j}}f(t)\,\mathrm dt\quad(x\in I_{n,j}).
\]
在每个分割原子上的积分保持，说明 \(A_nf=\mathbb E[f\mid\mathcal D_n]\)。
所有二进区间生成 Borel \(\sigma\)-代数：它们是 Borel 集；
反过来，每个相对开集都可写成包含于其中的二进区间的并，
因为包含某点的区间直径趋零，而且全部这样的区间只有可数多个。
故上一条定理给出
\[
 A_nf\longrightarrow f\quad\text{几乎处处且于 }L^1([0,1)).
\]
若 \(f\) 只是 Lebesgue 可测，由\cref{b1:thm:lebesgue-borel-representative}
先取一个与它几乎处处相等的 Borel 代表，
所得积分和结论不变。这里使用的是越来越细的平均，并不要求 \(f\) 处处连续。

\ProseParagraphBreak

例如 \(f(x)=x\) 时，\(A_nf\) 在每个区间上等于其中点，
所以 \(|A_nf(x)-x|\leq2^{-n-1}\)，甚至一致收敛。
对 \(f=\mathbf1_{[0,1/3)}\)，只有跨过 \(1/3\) 的那个区间需要取非平凡平均，
其余区间上已经准确恢复 \(f\)。跨界区间内的平均值在 \(1/3\) 与 \(2/3\) 之间交替，
因此 \(\|A_nf-f\|_\infty=2/3\) 对每个 \(n\geq0\) 成立，
尽管它仍几乎处处及于每个有限 \(L^p\) 空间收敛。
这是上一条推论中 \(L^\infty\) 结论不能补上的具体原因。

\subsection{非一致可积模型与失效现象}
\label{b1:subsec:310604}

一致可积性并非一个仅在证明中使用的技术补充，它决定能否在极限中保留条件平均关系。
对\eqref{b1:eq:dyadic-spike-martingale-31}，固定任意 \(K>0\)，再选 \(n\) 使 \(2^n>K\)，则
\[
 \mathbb E\bigl[M_n\cdot\mathbf1_{\{M_n>K\}}\bigr]=1.
\]
所以尾部积分的上确界始终为 \(1\)，并不趋零。
对每个有限时刻 \(N\)，前 \(N\) 项都由 \(M_N\) 闭合；
但随着 \(N\) 改变，这些闭合变量没有一个共同的可积终值可以替代。
若有这样的终值，\cref{b1:thm:closed-martingale-31}就会迫使全列一致可积，与上面的计算矛盾。

\ProseParagraphBreak

对非负鞅，还能用期望是否在极限中损失来判断。

\begin{proposition}[非负鞅的期望保持判据]
\label{b1:prop:nonnegative-martingale-ui-31}
设 \((\Omega,\mathcal F,\mathbb P)\) 为概率空间，\((M_n)_{n\geq0}\)
为其中适应于递增滤过的非负实鞅，\(M_\infty\) 为其几乎必然极限。
则 \(M_\infty\in L^1\)、\(\mathbb E M_\infty\leq\mathbb E M_0\)，并且
\[
 \{M_n:n\geq0\}\text{ 一致可积}
 \quad\Longleftrightarrow\quad
 \mathbb E M_\infty=\mathbb E M_0.
\]
\end{proposition}
\begin{proof}
非负鞅的一阶矩恒等于 \(\mathbb E M_0\)，故\cref{b1:cor:martingale-as-convergence-31}
保证存在可积的几乎必然极限，Fatou 引理给出期望不等式。
一致可积性推出 \(L^1\) 收敛，因而推出期望相等。
反过来，若期望相等，由非负性，
\[
 |M_n-M_\infty|=M_n+M_\infty-2\min(M_n,M_\infty).
\]
最后一项几乎必然趋于 \(M_\infty\)，并受可积变量 \(M_\infty\) 控制，
所以控制收敛给出
\[
 \mathbb E|M_n-M_\infty|
 \longrightarrow\mathbb E M_0-\mathbb E M_\infty=0.
\]
再用\cref{b1:thm:closed-martingale-31}即得一致可积性。
\end{proof}

这条判据的非负性不能删去。
将\eqref{b1:eq:dyadic-spike-martingale-31}的概率空间扩成
\([0,1)\times\{-1,1\}\)，配备乘积 Borel \(\sigma\)-代数和
长度测度与等概率符号分布的乘积测度。
取滤过 \(\mathcal D_n\otimes2^{\{-1,1\}}\)，即从时刻零就揭示符号坐标
\(\varepsilon\)，而第一个坐标仍只揭示 \(n\) 级二进信息。
过程 \(N_n(x,\varepsilon)=\varepsilon M_n(x)\) 仍是鞅：
在每个符号纤维的二进分块上，条件平均关系与原来相同。
它几乎必然趋于零，且每个 \(\mathbb E N_n=0=\mathbb E N_\infty\)，
但 \(\mathbb E|N_n|=1\)，所以不于 \(L^1\) 收敛，也不一致可积。
有符号的期望可能因正负部分抵消而保持，必须控制绝对值的积分质量。

\ProseParagraphBreak

本章中三种常用条件的后果列在\cref{b1:tab:martingale-convergence-31}。
表中假定过程始终为同一概率空间、同一递增滤过下的离散鞅；
孤立地把某个积分条件加在任意函数列上，并不能得到相同结论。

\begin{table}[htbp]
\centering
\caption{离散鞅的积分条件与收敛结论}
\label{b1:tab:martingale-convergence-31}
\begin{tabular}{p{0.34\textwidth}p{0.56\textwidth}}
\toprule
条件 & 保证的结论 \\
\midrule
\(\sup_n\mathbb E|M_n|<\infty\) & 几乎必然收敛，可积极限；不保证 \(L^1\) 收敛或终值表示。 \\
\(\{M_n\}\) 一致可积 & 几乎必然及 \(L^1\) 收敛，并有 \(M_n=\mathbb E[M_\infty\mid\mathcal F_n]\)。 \\
\(\sup_n\mathbb E|M_n|^p<\infty\)，\(1<p<\infty\) & 几乎必然及 \(L^p\) 收敛；最大变量 \(\sup_n|M_n|\) 属于 \(L^p\)。 \\
\bottomrule
\end{tabular}
\end{table}
